\documentclass{amsart}
% For dealing with references we use the comment environment
\usepackage{verbatim}
\newenvironment{reference}{\comment}{\endcomment}
%\newenvironment{reference}{}{}
\newenvironment{slogan}{\comment}{\endcomment}
\newenvironment{history}{\comment}{\endcomment}

% For commutative diagrams we use Xy-pic
\usepackage[all]{xy}

% We use 2cell for 2-commutative diagrams.
\xyoption{2cell}
\UseAllTwocells

% We use multicol for the list of chapters between chapters
\usepackage{multicol}

% This is generally recommended for better output
\usepackage{lmodern}
\usepackage[T1]{fontenc}

% For cross-file-references

% Package for hypertext links:
\usepackage{hyperref}

% For any local file, say "hello.tex" you want to link to please

% Theorem environments.
%
\theoremstyle{plain}
\newtheorem{theorem}[subsection]{Theorem}
\newtheorem{proposition}[subsection]{Proposition}
\newtheorem{lemma}[subsection]{Lemma}

\theoremstyle{definition}
\newtheorem{definition}[subsection]{Definition}
\newtheorem{example}[subsection]{Example}
\newtheorem{exercise}[subsection]{Exercise}
\newtheorem{situation}[subsection]{Situation}

\theoremstyle{remark}
\newtheorem{remark}[subsection]{Remark}
\newtheorem{remarks}[subsection]{Remarks}

\numberwithin{equation}{subsection}

% Macros
%
\def\lim{\mathop{\mathrm{lim}}\nolimits}
\def\colim{\mathop{\mathrm{colim}}\nolimits}
\def\Spec{\mathop{\mathrm{Spec}}}
\def\Hom{\mathop{\mathrm{Hom}}\nolimits}
\def\Ext{\mathop{\mathrm{Ext}}\nolimits}
\def\SheafHom{\mathop{\mathcal{H}\!\mathit{om}}\nolimits}
\def\SheafExt{\mathop{\mathcal{E}\!\mathit{xt}}\nolimits}
\def\Sch{\mathit{Sch}}
\def\Mor{\mathop{\mathrm{Mor}}\nolimits}
\def\Ob{\mathop{\mathrm{Ob}}\nolimits}
\def\Sh{\mathop{\mathit{Sh}}\nolimits}
\def\NL{\mathop{N\!L}\nolimits}
\def\CH{\mathop{\mathrm{CH}}\nolimits}
\def\proetale{{pro\text{-}\acute{e}tale}}
\def\etale{{\acute{e}tale}}
\def\QCoh{\mathit{QCoh}}
\def\Ker{\mathop{\mathrm{Ker}}}
\def\Im{\mathop{\mathrm{Im}}}
\def\Coker{\mathop{\mathrm{Coker}}}
\def\Coim{\mathop{\mathrm{Coim}}}

% Boxtimes
%
\DeclareMathSymbol{\boxtimes}{\mathbin}{AMSa}{"02}

%
% Macros for moduli stacks/spaces
%
\def\QCohstack{\mathcal{QC}\!\mathit{oh}}
\def\Cohstack{\mathcal{C}\!\mathit{oh}}
\def\Spacesstack{\mathcal{S}\!\mathit{paces}}
\def\Quotfunctor{\mathrm{Quot}}
\def\Hilbfunctor{\mathrm{Hilb}}
\def\Curvesstack{\mathcal{C}\!\mathit{urves}}
\def\Polarizedstack{\mathcal{P}\!\mathit{olarized}}
\def\Complexesstack{\mathcal{C}\!\mathit{omplexes}}
% \Pic is the operator that assigns to X its picard group, usage \Pic(X)
% \Picardstack_{X/B} denotes the Picard stack of X over B
% \Picardfunctor_{X/B} denotes the Picard functor of X over B
\def\Pic{\mathop{\mathrm{Pic}}\nolimits}
\def\Picardstack{\mathcal{P}\!\mathit{ic}}
\def\Picardfunctor{\mathrm{Pic}}
\def\Deformationcategory{\mathcal{D}\!\mathit{ef}}

\setlength{\emergencystretch}{3em}
\begin{document}
\title[Stacks Project, Tag 09T6]{379 --- Local derived pushforward compatibility implies global quasi-coherent derived-category equivalence}
\maketitle
\noindent Copyright (C) 2005--2025 Johan de Jong. Stacks Project authors. Source: \href{https://stacks.math.columbia.edu/tag/09T6}{Tag 09T6}.

\noindent Editorial difficulty note (not author text). Difficulty H2: The unit and counit of the derived coherator must be proved invertible for every unbounded object. The proof inducts on the number of affine opens meeting the closed support, constructs localization triangles and uses coherator-pushforward compatibility to reduce both maps to the affine case..

\noindent Editorial provenance note (not author text). History: excerpt from revision \texttt{a0444\allowbreak 6e57e\allowbreak c1fbc\allowbreak 252a8\allowbreak 71afc\allowbreak ec775\allowbreak 2fb28\allowbreak 07b14}, perfect.tex, lines 1484--1581. Reference presentation was adapted; original-author context and core support blocks were retained or added. The mathematical wording is unchanged. Licensed under GNU FDL 1.2 or later; no invariant sections or cover texts. See ../licenses/COPYING. Context blocks reproduce author definitions and supporting results; they are not additional counted problems.

\begin{definition}
\label{definition-supported-on}
Let $X$ be a scheme. Let $E$ be an object of $D(\mathcal{O}_X)$.
Let $T \subset X$ be a closed subset.
We say $E$ is {\it supported on $T$} if the
cohomology sheaves $H^i(E)$ are supported on $T$.
\end{definition}

\noindent
In this section we discuss the relationship between quasi-coherent
modules and all modules on a scheme $X$. A reference is
\href{https://stacks.math.columbia.edu/bibliography}{Bibliography: TT (Appendix B)}. By the discussion in
Schemes, Section \href{https://stacks.math.columbia.edu/tag/01LA}{section quasi coherent (Tag 01LA)}
the embedding
$\QCoh(\mathcal{O}_X) \subset \textit{Mod}(\mathcal{O}_X)$
exhibits $\QCoh(\mathcal{O}_X)$ as a weak Serre subcategory of
the category of $\mathcal{O}_X$-modules. Denote
$$
D_\QCoh(\mathcal{O}_X) \subset D(\mathcal{O}_X)
$$
the subcategory of complexes whose cohomology sheaves are quasi-coherent, see
Derived Categories, Section \href{https://stacks.math.columbia.edu/tag/06UP}{section triangulated sub (Tag 06UP)}.
Thus we obtain a canonical functor
\begin{equation}
\label{equation-compare}
D(\QCoh(\mathcal{O}_X))
\longrightarrow
D_\QCoh(\mathcal{O}_X)
\end{equation}
see Derived Categories, Equation (\href{https://stacks.math.columbia.edu/tag/06UR}{equation compare (Tag 06UR)}).

\section{The coherator}
\label{section-coherator}

\noindent
Let $X$ be a scheme. The {\it coherator} is a functor
$$
Q_X :
\textit{Mod}(\mathcal{O}_X)
\longrightarrow
\QCoh(\mathcal{O}_X)
$$
which is right adjoint to the inclusion functor
$\QCoh(\mathcal{O}_X) \to \textit{Mod}(\mathcal{O}_X)$.
It exists for any scheme $X$ and moreover the adjunction mapping
$Q_X(\mathcal{F}) \to \mathcal{F}$ is an isomorphism for every
quasi-coherent module $\mathcal{F}$, see
Properties, Proposition \href{https://stacks.math.columbia.edu/tag/077P}{proposition coherator (Tag 077P)}.
Since $Q_X$ is left exact (as a right adjoint) we can consider its
right derived extension
$$
RQ_X :
D(\mathcal{O}_X)
\longrightarrow
D(\QCoh(\mathcal{O}_X)).
$$
Since $Q_X$ is right adjoint to the inclusion functor
$\QCoh(\mathcal{O}_X) \to \textit{Mod}(\mathcal{O}_X)$
we see that $RQ_X$ is right adjoint to the canonical functor
$D(\QCoh(\mathcal{O}_X)) \to D(\mathcal{O}_X)$ by
Derived Categories, Lemma \href{https://stacks.math.columbia.edu/tag/09T5}{lemma derived adjoint functors (Tag 09T5)}.

\medskip\noindent
In this section we will study the functor $RQ_X$. In
Section \href{https://stacks.math.columbia.edu/tag/0CQZ}{section better coherator (Tag 0CQZ)}
we will study the (closely related) right adjoint to the inclusion functor
$D_\QCoh(\mathcal{O}_X) \to D(\mathcal{O}_X)$ (when it exists).



\begin{lemma}
\label{lemma-argument-proves}
Let $X$ be a quasi-compact and quasi-separated scheme. Suppose that
for every affine open $U \subset X$ the right derived functor
$$
\Phi : D(\QCoh(\mathcal{O}_U)) \to D(\QCoh(\mathcal{O}_X))
$$
of the left exact functor
$j_* : \QCoh(\mathcal{O}_U) \to \QCoh(\mathcal{O}_X)$
fits into a commutative diagram
$$
\xymatrix{
D(\QCoh(\mathcal{O}_U)) \ar[d]_\Phi \ar[r]_{i_U} &
D_\QCoh(\mathcal{O}_U) \ar[d]^{Rj_*} \\
D(\QCoh(\mathcal{O}_X)) \ar[r]^{i_X} &
D_\QCoh(\mathcal{O}_X)
}
$$
Then the functor (\ref{equation-compare})
$$
D(\QCoh(\mathcal{O}_X))
\longrightarrow
D_\QCoh(\mathcal{O}_X)
$$
is an equivalence with quasi-inverse given by $RQ_X$.
\end{lemma}

\begin{proof}
Let $E$ be an object of $D_\QCoh(\mathcal{O}_X)$ and
let $A$ be an object of $D(\QCoh(\mathcal{O}_X))$.
We have to show that the adjunction maps
$$
A \to RQ_X(i_X(A))
\quad\text{and}\quad
i_X(RQ_X(E)) \to E
$$
are isomorphisms. Consider the hypothesis
$H_n$: the adjunction maps above are isomorphisms
whenever $E$ and $i_X(A)$ are supported
(Definition \ref{definition-supported-on})
on a closed subset of $X$ which
is contained in the union of $n$ affine opens of $X$.
We will prove $H_n$ by induction on $n$.

\medskip\noindent
Base case: $n = 0$. In this case $E = 0$, hence the map
$i_X(RQ_X(E)) \to E$ is an isomorphism. Similarly $i_X(A) = 0$.
Thus the cohomology sheaves of $i_X(A)$ are zero. Since the inclusion
functor $\QCoh(\mathcal{O}_X) \to \textit{Mod}(\mathcal{O}_X)$
is fully faithful and exact, we conclude that the cohomology
objects of $A$ are zero, i.e., $A = 0$ and
$A \to RQ_X(i_X(A))$ is an isomorphism as well.

\medskip\noindent
Induction step. Suppose that $E$ and $i_X(A)$ are supported on a
closed subset $T$ of $X$ contained in $U_1 \cup \ldots \cup U_n$
with $U_i \subset X$ affine open. Set $U = U_n$.
Consider the distinguished triangles
$$
A \to \Phi(A|_U) \to A' \to A[1]
\quad\text{and}\quad
E \to Rj_*(E|_U) \to E' \to E[1]
$$
where $\Phi$ is as in the statement of the lemma.
Note that $E \to Rj_*(E|_U)$ is a quasi-isomorphism over $U = U_n$.
Since $i_X \circ \Phi = Rj_* \circ i_U$ by assumption
and since $i_X(A)|_U = i_U(A|_U)$
we see that $i_X(A) \to i_X(\Phi(A|_U))$ is a quasi-isomorphism over $U$.
Hence $i_X(A')$ and $E'$ are supported on the closed
subset $T \setminus U$ of $X$ which is contained in
$U_1 \cup \ldots \cup U_{n - 1}$.
By induction hypothesis the statement is true for $A'$ and $E'$. By
Derived Categories, Lemma \href{https://stacks.math.columbia.edu/tag/014A}{lemma third isomorphism triangle (Tag 014A)}
it suffices to prove the maps
$$
\Phi(A|_U) \to RQ_X(i_X(\Phi(A|_U)))
\quad\text{and}\quad
i_X(RQ_X(Rj_*E|_U)) \to Rj_*(E|_U)
$$
are isomorphisms. By assumption and by
Lemma \href{https://stacks.math.columbia.edu/tag/08D8}{lemma flat pushforward coherator (Tag 08D8)}
(the inclusion morphism $j : U \to X$ is flat, quasi-compact, and
quasi-separated) we have
$$
RQ_X(i_X(\Phi(A|_U))) = RQ_X(Rj_*(i_U(A|_U))) = \Phi(RQ_U(i_U(A|_U)))
$$
and
$$
i_X(RQ_X(Rj_*(E|_U))) = i_X(\Phi(RQ_U(E|_U))) = Rj_*(i_U(RQ_U(E|_U)))
$$
Finally, the maps
$$
A|_U \to RQ_U(i_U(A|_U))
\quad\text{and}\quad
i_U(RQ_U(E|_U)) \to E|_U
$$
are isomorphisms by Lemma \href{https://stacks.math.columbia.edu/tag/08D9}{lemma affine coherator (Tag 08D9)}. The result follows.
\end{proof}
\end{document}
